\section{Conclusion}
\label{sec:conclusion}

BAR distributes a nominal TD3+BC actor coefficient across a persistent actor chain, couples Bellman targets to the first branch, and deploys the last. Across a complete fixed-budget phase diagram, deeper proximal configurations substantially reduce large-budget failures, with the strongest tested configuration raising the high-budget mean from $25.97$ to $63.61$. Dataset-cluster analysis supports a broad rescue relative to one-step TD3+BC, while incremental collapse-risk evidence from $K=3$ to $K=4$ is weaker. Re-audited diagnostics show consistent contraction of sample-anchored target-action movement and severe critic miscalibration in screened failures. These findings establish a strong descriptive bundle effect. The proximal analysis locally motivates the construction but does not certify the implemented one-Adam-step actor chain or isolate its causal mechanism.
